\subsection{Кто в сети}
\label{sec:swarm-units}

Цель раздела: подключиться к рою, узнать свой идентификатор и список соседей и корректно закрыть соединение.

\begin{codefile}[python]{python-onboard/swarm/swarm_01_who_is_here.py}
\end{codefile}

\paragraph{Что происходит.} \verb|PioneerSwarmUnit()| подключает программу к рою, \verb|swarm.id| --- свой \texttt{machine\_id}. Свой тип известен только после первого собственного Heartbeat, поэтому цикл \verb|while swarm.unit_type is None| ждёт его: до этого свойство равно \verb|None|. Затем пять раз в секунду печатается список соседей. Он не содержит самого дрона и меняется на лету: новый юнит появляется после своего первого Heartbeat (в симуляторе --- примерно через полсекунды), пропавший --- через три периода тишины. В конце \verb|swarm.close()| освобождает подписки. Остальные дроны сцены отвечать не обязаны: \texttt{swarm-server} каждого борта объявляет его сам.

\begin{taskbox}[Кто в рое]
Добавьте в сцену два дрона Мини~2 и запустите программу на первом. Проверьте, что программа печатает свой \texttt{machine\_id}, тип и список соседей, и что себя в списке нет. Затем дополните её: для каждого соседа печатайте отдельную строку «Сосед: \ldots».
\expected{у дрона 1 --- «Я: pioneer1» и «Мой тип: SwarmUnitType.AERO (воздушный)», соседи \texttt{pioneer2} и \texttt{pioneer3} в каждой из пяти строк списка и по строке «Сосед:» на каждого.}
\mission{Мини~2 SDK2}{57}
\end{taskbox}

\subsection{Сообщения: ящик и опрос}
\label{sec:swarm-messages}

Цель раздела: отправлять и принимать текстовые сообщения и строить на них простой протокол.

Сообщение --- строка. Её отправляют в ящик конкретного юнита:

\begin{codebox}[Отправка и приём]
\begin{verbatim}
# True или False
ok = swarm.send_message_to_unit('pioneer2', 'ping 1')
# (sender, text) или None
message = swarm.get_next_message()
\end{verbatim}
\end{codebox}

\verb|send_message_to_unit| возвращает \verb|True|, если адресат найден и сообщение ушло, и \verb|False|, если юнита в списке нет --- в том числе когда вы пишете самому себе. Содержимое протокол не проверяет, но и доставку не подтверждает: \verb|True| означает «ушло», а не «прочитано». Приём устроен иначе, чем у обычных подписок: \verb|get_next_message()| \emph{не ждёт} и возвращает \verb|None|, когда ящик пуст; ждать сообщение нужно самому --- опросом с паузой и тайм-аутом.

\begin{rulebox}[Опрос ящика --- с паузой и пределом]
Цикл ожидания сообщения всегда имеет паузу (\verb|time.sleep(0.1)|), иначе он занимает все ресурсы процессора, и предел (счётчик времени или число попыток), иначе программа зависнет навсегда, если собеседник не ответит. При выходе по пределу программа решает сама: повторить, отказаться, сообщить оператору.
\end{rulebox}

\begin{infobox}[Свойства ящика в симуляторе]
Доставка мгновенная, порядок сохраняется. Ящик принадлежит юниту, а не программе: письмо соседу, у которого нет программы, ждёт, пока она появится. Ящик ограничен 100 сообщениями (старые вытесняются), очищается, когда закрывается последний объект юнита (\verb|close()| или конец программы), и пропадает вместе с юнитом. Несколько объектов на одном дроне читают из одного ящика. На настоящей сети эти свойства --- допущения: доставку и порядок там никто не гарантирует.
\end{infobox}

\subsubsection*{Ping-pong}

Классическая проверка связи: один дрон отправляет «ping N» и ждёт «pong N», другой на каждое «ping» отвечает «pong». Два файла --- для двух дронов; первым запускают отвечающего.

\begin{codefile}[python]{python-onboard/swarm/swarm_02_ping_pong_pong.py}
\end{codefile}

\paragraph{Что происходит (pong).} Дрон ждёт сообщение в ящике. Пустой ящик --- пауза 0.1~с и накопление счётчика тишины \verb|silence|; пришло сообщение --- счётчик обнуляется, из адресата и текста выделяется ответ: \verb|'pong' + text[4:]| заменяет первые четыре буквы «ping» на «pong», оставляя номер. Ответ уходит тому, кто написал: \verb|sender|. Десять секунд тишины --- программа закрывает соединение.

\begin{codefile}[python]{python-onboard/swarm/swarm_02_ping_pong_ping.py}
\end{codefile}

\paragraph{Что происходит (ping).} Первый сосед из списка --- собеседник (\verb|partner|). Пять раз: отправка «ping N», затем ожидание ответа до трёх секунд (\verb|waited|) опросом с паузой. Если сосед не отвечает, строка «Ответа нет» --- и цикл идёт дальше. Если соседей нет вовсе, программа подсказывает добавить дрон.

\begin{taskbox}[Ping-pong]
Запустите \texttt{swarm\_02\_ping\_pong\_pong.py} на втором дроне, затем \texttt{swarm\_02\_ping\_pong\_ping.py} на первом. Измените протокол: в ответ добавляйте время ожидания и свой \texttt{machine\_id} (\texttt{pong 3 от pioneer2}). Что произойдёт, если запустить «ping» раньше «pong»? Объясните по правилу об опросе ящика.
\expected{первый дрон печатает пять строк «Ответ от pioneer2 : pong N»; второй --- пять строк «Получил от pioneer1 : ping N»; если «pong» запущен позже, письма ждут в его ящике до старта программы (ящик принадлежит юниту).}
\mission{Мини~2 SDK2}{58, 59}
\end{taskbox}

\subsubsection*{Рассылка всем соседям}

В рое нет «сообщения всем»: письмо уходит каждому юниту отдельно, а общая рассылка --- цикл по списку.

\begin{codefile}[python]{python-onboard/swarm/swarm_03_broadcast_receiver.py}
\end{codefile}

\begin{codefile}[python]{python-onboard/swarm/swarm_03_broadcast_sender.py}
\end{codefile}

\paragraph{Что происходит.} Получатель (его запускают на втором и третьем дронах первым) восемь секунд опрашивает ящик и печатает всё пришедшее со счётчиком. Отправитель выжидает две секунды, чтобы получатели успели запуститься, затем перебирает соседей и считает, кому письмо ушло. Вызов для несуществующего юнита (\texttt{pioneer99}) возвращает \verb|False|: эту проверку используйте всегда, вместо того чтобы считать, что письмо дошло.

\begin{taskbox}[Рассылка и подтверждение]
Расширьте рассылку: отправитель отправляет «кто здесь?» всем соседям, а каждый получатель отвечает «здесь $\langle$id$\rangle$». Отправитель печатает «Ответили: N из M» и строку для каждого ответившего.
\expected{при двух получателях --- «Ответили: 2 из 2» и по строке «pioneerN - здесь pioneerN» на каждого; если один получатель не запущен, число ответивших меньше числа соседей.}
\mission{Мини~2 SDK2}{60}
\end{taskbox}

\paragraph{Протокол из строк.} Сообщение --- просто строка, поэтому формат придумываете вы. Для команд подходят короткие слова (\verb|'взлёт'|, \verb|'отбой'|); для данных --- слово и значения через пробел (\verb|'НАЙДЕН 2.00 1.50'|), которые разбирают \verb|text.split()| и \verb|float|. Правила хорошего протокола: одно сообщение --- одна мысль; команда в начале, данные после; при разборе проверять число частей и допускать мусор (чужое сообщение не должно ронять программу). Примеры дальше (синхронный взлёт, световое шоу, поиск и спасение) --- именно такие протоколы.
